\chapter*{Capítulo III}
\addtocontents{toc}{\protect\contentsline{chapter}{Capítulo III}{\thepage}{}}
\setcounter{chapter}{3} % Set the chapter counter to ensure proper numbering
\setcounter{section}{0} % Set the chapter counter to ensure proper numbering
\fancyhf{}
\rhead{Capítulo III}
\lhead{Informe de práctica Inicial en Air Logistics Corp.}
\fancyfoot[R]{\thepage} 
\section{Conclusiones}
La experiencia de trabajar en una empresa real fue invaluable. Poder presenciar personalmente los problemas y la complejidad de los procesos de una empresa significaron una buena base para así en las asignaturas posteriores tener una manera de tener casos reales para poder proponer soluciones utilizando los métodos aprendidos posteriormente. 
\subsection{Resultados obtenidos.}
Al finalizar el periodo de práctica inicial, el personal de la empresa Air Logistics quedó satisfecho con los sistemas desarrollados durante la práctica. Estos sistemas han de ahorrar una significativa cantidad de tiempo y dinero a esta empresa.
\subsection{Relación con la carrera.}
Las asignaturas más relacionadas a las actividades de esta práctica fueron Programación Orientada a Objetos y Sistemas Corporativos de Gestión. A pesar que estas asignaturas tienden a ser más informáticas, se reconoce que el manejo de la información y la automatización son críticas en la Ingeniería industrial.
\subsection{Reflexión.}
Fue desafiante el tener que aprender la mayoría de las cosas que serían utilizadas en el desarrollo de los sistemas. Hubieron varias veces que surgió el desánimo, sin embargo, hubo una gran satisfacción al saber que lo que se realizó en la empresa tuvo un impacto positivo.
